% GENERATED FILE. Edit canonical lessons, then run npm run generate:books.
% canonical-source-hash: fnv1a64:5b620c65dc28b967
% canonical-lessons: ES-C754-tejer, ES-C754-pesca, ES-C754-ceramica, ES-C754-jardineria, ES-C754-coser

\chapter{Aficiones tranquilas}
\label{ch:es-aficiones-tranquilas}
\hlchaptermodality{754}

\begin{chapteropening}
\textbf{By the end of this chapter:} \emph{I can talk about quiet hobbies: fishing, gardening, sewing, knitting and pottery.}

Complete the last lesson of chapter 754: I can talk about quiet hobbies: fishing, gardening, sewing, knitting and pottery.
\end{chapteropening}

\section[tejer]{tejer — to knit, to weave}

\label{lesson:ES-C754-tejer}



\begin{quote}
Before the new one: say the Spanish for rhythm, then the Spanish for to compose.
\end{quote}

\subsection*{You'll want to know: tejer}
\textbf{tejer} — \textquotedblleft{}to knit, to weave\textquotedblright{}.

\begin{quote}
Mi abuela teje bufandas.
\end{quote}

\begin{quote}
\emph{My grandmother knits scarves.}
\end{quote}

Spain also says \emph{hacer punto}.

\begin{cousinweb}
\emph{tejer} is inherited from Latin \emph{texere} \textquotedblleft{}to weave\textquotedblright{} (EtymDB); a text is words woven together.
\end{cousinweb}

\begin{grammarlens}[title={What you've built}]
One more word for telling, explaining and describing at B1.
\end{grammarlens}

\subsection*{Your turn}
\begin{itemize}
\raggedright
  \item \emph{Say it:} \emph{tejer}
  \item \emph{Say it:} \emph{tejer}, once more
  \item \emph{Say it:} a sentence of your own with \emph{tejer}
  \item \emph{Recall:} say the Spanish for rhythm, then the Spanish for to compose, then say \emph{tejer} again
\end{itemize}

\begin{tcolorbox}[breakable,colback=teal!4,colframe=teal!35!black,title={Before you move on}]
What does tejer mean? (\textquotedblleft{}To knit, to weave\textquotedblright{}.) Say it once more.
\end{tcolorbox}
\section[la pesca]{la pesca — fishing}

\label{lesson:ES-C754-pesca}



\begin{quote}
Before the new one: say the Spanish for to compose, then the Spanish for to knit.
\end{quote}

\subsection*{You'll want to know: la pesca}
\textbf{la pesca} — \textquotedblleft{}fishing\textquotedblright{}.

\begin{quote}
La pesca es su gran afición.
\end{quote}

\begin{quote}
\emph{Fishing is his great hobby.}
\end{quote}

\begin{cousinweb}
\emph{pesca} is built on \emph{pescar} (EtymDB), which is inherited from Latin \emph{piscārī} \textquotedblleft{}to fish\textquotedblright{}.
\end{cousinweb}

\begin{grammarlens}[title={What you've built}]
One more word for telling, explaining and describing at B1.
\end{grammarlens}

\subsection*{Your turn}
\begin{itemize}
\raggedright
  \item \emph{Say it:} \emph{la pesca}
  \item \emph{Say it:} \emph{la pesca}, once more
  \item \emph{Say it:} a sentence of your own with \emph{la pesca}
  \item \emph{Recall:} say the Spanish for to compose, then the Spanish for to knit, then say \emph{la pesca} again
\end{itemize}

\begin{tcolorbox}[breakable,colback=teal!4,colframe=teal!35!black,title={Before you move on}]
What does la pesca mean? (\textquotedblleft{}Fishing\textquotedblright{}.) Say it once more.
\end{tcolorbox}
\section[la cerámica]{la cerámica — pottery, ceramics}

\label{lesson:ES-C754-ceramica}



\begin{quote}
Before the new one: say the Spanish for to knit, then the Spanish for fishing.
\end{quote}

\subsection*{You'll want to know: la cerámica}
\textbf{la cerámica} — \textquotedblleft{}pottery, ceramics\textquotedblright{}.

\begin{quote}
Hago un curso de cerámica.
\end{quote}

\begin{quote}
\emph{I am doing a pottery course.}
\end{quote}

\begin{cousinweb}
\emph{cerámica} comes from Greek \emph{keramik\'{\={e}}} \textquotedblleft{}the potter's art\textquotedblright{}, built on \emph{kéramos} \textquotedblleft{}potter's clay\textquotedblright{}; it probably came in through French \emph{céramique}.
\end{cousinweb}

\begin{grammarlens}[title={What you've built}]
One more word for telling, explaining and describing at B1.
\end{grammarlens}

\subsection*{Your turn}
\begin{itemize}
\raggedright
  \item \emph{Say it:} \emph{la cerámica}
  \item \emph{Say it:} \emph{la cerámica}, once more
  \item \emph{Say it:} a sentence of your own with \emph{la cerámica}
  \item \emph{Recall:} say the Spanish for to knit, then the Spanish for fishing, then say \emph{la cerámica} again
\end{itemize}

\begin{tcolorbox}[breakable,colback=teal!4,colframe=teal!35!black,title={Before you move on}]
What does la cerámica mean? (\textquotedblleft{}Pottery, ceramics\textquotedblright{}.) Say it once more.
\end{tcolorbox}
\section[la jardinería]{la jardinería — gardening}

\label{lesson:ES-C754-jardineria}



\begin{quote}
Before the new one: say the Spanish for fishing, then the Spanish for pottery.
\end{quote}

\subsection*{You'll want to know: la jardinería}
\textbf{la jardinería} — \textquotedblleft{}gardening\textquotedblright{}.

\begin{quote}
Mi abuelo dedica la mañana a la jardinería.
\end{quote}

\begin{quote}
\emph{My grandfather spends the morning gardening.}
\end{quote}

\begin{cousinweb}
\emph{jardinería} is built on \emph{jardín}, which is usually traced to French \emph{jardin}, from Frankish \emph{gardo} \textquotedblleft{}enclosure\textquotedblright{}.
\end{cousinweb}

\begin{grammarlens}[title={What you've built}]
One more word for telling, explaining and describing at B1.
\end{grammarlens}

\subsection*{Your turn}
\begin{itemize}
\raggedright
  \item \emph{Say it:} \emph{la jardinería}
  \item \emph{Say it:} \emph{la jardinería}, once more
  \item \emph{Say it:} a sentence of your own with \emph{la jardinería}
  \item \emph{Recall:} say the Spanish for fishing, then the Spanish for pottery, then say \emph{la jardinería} again
\end{itemize}

\begin{tcolorbox}[breakable,colback=teal!4,colframe=teal!35!black,title={Before you move on}]
What does la jardinería mean? (\textquotedblleft{}Gardening\textquotedblright{}.) Say it once more.
\end{tcolorbox}
\section[coser]{coser — to sew}

\label{lesson:ES-C754-coser}



\begin{quote}
Before the new one: say the Spanish for pottery, then the Spanish for gardening.
\end{quote}

\subsection*{You'll want to know: coser}
\textbf{coser} — \textquotedblleft{}to sew\textquotedblright{}.

\begin{quote}
Mi madre me enseñó a coser.
\end{quote}

\begin{quote}
\emph{My mother taught me to sew.}
\end{quote}

\emph{Ser coser y cantar} — to be a piece of cake.

\begin{cousinweb}
\emph{coser} is inherited from Latin \emph{cōnsuere} \textquotedblleft{}to sew together\textquotedblright{}, through spoken \emph{cosēre} (EtymDB).
\end{cousinweb}

\begin{grammarlens}[title={What you've built}]
One more word for telling, explaining and describing at B1.
\end{grammarlens}

\subsection*{Your turn}
\begin{itemize}
\raggedright
  \item \emph{Say it:} \emph{coser}
  \item \emph{Say it:} \emph{coser}, once more
  \item \emph{Say it:} a sentence of your own with \emph{coser}
  \item \emph{Recall:} say the Spanish for pottery, then the Spanish for gardening, then say \emph{coser} again
\end{itemize}

\begin{tcolorbox}[breakable,colback=teal!4,colframe=teal!35!black,title={Before you move on}]
What does coser mean? (\textquotedblleft{}To sew\textquotedblright{}.) Say it once more.
\end{tcolorbox}
